\documentclass[10pt,a4paper]{amsart}
\usepackage[margin=19mm]{geometry}
\usepackage{amsmath,amssymb,mathtools}
\usepackage[scr=rsfs,cal=euler]{mathalfa}
\usepackage{xfrac}
\usepackage[unicode,hidelinks,bookmarks=false]{hyperref}
\newcommand{\ie}{i.e.\ }
\newcommand{\cf}{cf.\ }
\newcommand{\resp}{respectively\ }
\newcommand{\Subsection}{\subsection}
\providecommand{\nobreakdash}{\nobreak\hbox{-}}
\setlength{\emergencystretch}{2em}
\multlinegap0pt
\newtheorem{thm}{Theorem}[section]
\newtheorem{lem}[thm]{Lemma}
\newtheorem{prop}[thm]{Proposition}
\newtheorem{cor}[thm]{Corollary}
\newtheorem{defn}[thm]{Definition}
\newtheorem{example}[thm]{Example}
\newtheorem{remark}[thm]{Remark}
\setlength{\parskip}{4pt}
\numberwithin{equation}{section}
\begin{document}
\title[Higher Schwarzians of Elliptic Double Covers and Eisenstein-]{Higher Schwarzians of Elliptic Double Covers and Eisenstein--Kronecker Functions}
\author{Hicham Saber; Abdellah Sebbar}
\date{}
\maketitle
\noindent\emph{Source and licence.} Higher Schwarzians of Elliptic Double Covers and Eisenstein--Kronecker Functions. Version arXiv:2609.16444v1 (CC BY 4.0). This material is distributed under the Creative Commons Attribution 4.0 International licence, http://creativecommons.org/licenses/by/4.0/, as recorded in the arXiv OAI licence metadata for this exact version and shown on the abstract page https://arxiv.org/abs/2609.16444v1. Full rights notice: licenses/arxiv-2609.16444-CC-BY-4.0.txt.\par\smallskip
\noindent\textit{Editorial scope.} Counted unit: the original \texttt{thm} environment \texttt{thm:general-Hasse} at source lines 1469--1491 together with its complete proof at lines 1493--1515; the delivered document keeps source lines 1337--1516 so that every referenced label and every citation resolves inside it. Native editable source is the paper's own concatenated TeX; the AMS wrapper is editorial formatting only. No publisher class, upstream script or unrelated artwork is redistributed.
\section{Bernoulli reduction and the Hasse invariant}
\label{sec:hasse-reduction}

The reduction in order $p-1$ has a differential origin independent of
the degree-two hypothesis. We prove the local identity first and then
specialize it to invariant derivations on elliptic curves. The
universal double cover is a global corollary. Throughout this section
$p\geq5$ is prime and the Bernoulli numbers are normalized by $B_2=1/6$.

\subsection{A general differential reduction formula}

Let $B$ be a $p$-torsion-free $\mathbb Z_{(p)}$-algebra, let $\partial$
be a $\mathbb Z_{(p)}$-derivation of $B$, and let $f\in B$ with
$\partial f\in B^\times$. In $B[1/p]((t))$, define
\begin{equation}\label{eq:differential-Aharonov}
\frac{\partial f}{\displaystyle\sum_{j\geq1}
 (\partial^j f)t^j/j!}
=\frac1t-\sum_{n\geq1}S_n^\partial[f]t^{n-1}.
\end{equation}
This is the usual Aharonov expansion when $\partial=d/dz$. A bar
will denote reduction modulo $p$.

\begin{prop}\label{prop:Bernoulli-reduction}
Under these hypotheses, the normalized expression
\[
\mathcal S_{p-1}^\partial[f]
:=-\frac{(p-1)!}{B_{p-1}}S_{p-1}^\partial[f]
\]
belongs to $B$, and
\begin{equation}\label{eq:universal-Bernoulli-reduction}
\overline{\mathcal S_{p-1}^\partial[f]}
=\frac{\bar\partial^{\,p}\bar f}{\bar\partial\bar f}
\quad\text{in } B/pB.
\end{equation}
More precisely,
\begin{equation}\label{eq:leading-p-derivative}
S_{p-1}^\partial[f]
=\frac{\partial^p f}{p!\,\partial f}+R_p[f],
\qquad R_p[f]\in B.
\end{equation}
\end{prop}

\begin{proof}
Put
\[
b_j=\frac{\partial^{j+1}f}{(j+1)!\,\partial f}.
\]
The left side of \eqref{eq:differential-Aharonov} is
$t^{-1}(1+\sum_{j\geq1}b_jt^j)^{-1}$. In the coefficient of $t^{p-1}$
of the inverse series, the only term containing $b_{p-1}$ is
$-b_{p-1}$; every other term is an integral polynomial in
$b_1,\ldots,b_{p-2}$. For $j\leq p-2$, the denominator $(j+1)!$ is a
unit in $\mathbb Z_{(p)}$. Comparing this coefficient with
$-S_{p-1}^\partial[f]$ proves \eqref{eq:leading-p-derivative}.

The von Staudt--Clausen theorem gives
\begin{equation}\label{eq:Bernoulli-p-unit}
pB_{p-1}\in\mathbb Z_{(p)}^\times,
\qquad pB_{p-1}\equiv-1\pmod p
\end{equation}
\cite[pp.~233--237]{ireland-rosen}. Multiplying
\eqref{eq:leading-p-derivative} by $-(p-1)!/B_{p-1}$ gives
\[
\mathcal S_{p-1}^\partial[f]
=-\frac1{pB_{p-1}}\frac{\partial^p f}{\partial f}
 -\frac{(p-1)!}{B_{p-1}}R_p[f].
\]
Both terms belong to $B$. The coefficient of $R_p[f]$ lies in
$p\mathbb Z_{(p)}$, while $-1/(pB_{p-1})$ reduces to $1$. This proves
\eqref{eq:universal-Bernoulli-reduction}.
\end{proof}

This proof shows that the normalized expression is a universal
$\mathbb Z_{(p)}$-polynomial in the derivatives through order $p$, with
powers of $\partial f$ inverted. Thus its reduction depends only on
the reduced differential algebra and not on an integral lift. It is
this integral expression which is reduced; no Taylor expansion with
$1/p!$ as a coefficient is asserted in characteristic $p$.

\begin{remark}\label{rem:one-dimensional-reduction}
Map independence already holds for rational derivations on a curve.
Let $C$ be a smooth geometrically integral curve over a perfect field
$k$ of characteristic $p$, put $K=k(C)$, and let
$D\in\operatorname{Der}_k(K)$ be nonzero. Since
$\Omega^1_{K/k}$ has dimension one over $K$, so does
$\operatorname{Der}_k(K)$. The iterated Leibniz rule makes $D^p$ a
$k$-derivation, because all the intermediate binomial coefficients
vanish in characteristic $p$. Consequently there is a unique
$a_D\in K$ such that
\[
D^p=a_DD,\qquad \frac{D^pf}{Df}=a_D\quad\text{whenever }Df\neq0.
\]
The characteristic-$p$ evaluation of the universal polynomial in
Proposition~\ref{prop:Bernoulli-reduction} is therefore $a_D$ for
every such $f$. In general $a_D$ is a rational function on $C$.
For an elliptic invariant derivation, the coefficient is pulled back
from the base and is identified below with the Hasse invariant.
\end{remark}

\subsection{Independence of the rational function on an elliptic curve}

Let $E$ be an elliptic curve over an $\mathbb F_p$-algebra $R$, let
$\omega$ generate its invariant differentials, and let $\partial$ be
the invariant derivation dual to $\omega$. Write
\[
A(E)=a_\omega\,\omega^{\otimes(p-1)}
\]
for the Hasse invariant. Thus $a_\omega=A(E,\omega)$ is its coefficient
in this frame.

\begin{lem}[Restricted power of the invariant derivation]
\label{lem:restricted-Hasse}
One has
\begin{equation}\label{eq:restricted-Hasse}
\partial^p=a_\omega\partial.
\end{equation}
\end{lem}

\begin{proof}
In characteristic $p$, the iterated Leibniz rule shows that $\partial^p$
is a derivation, since $\binom pj=0$ for $0<j<p$. It is again
translation-invariant and hence is a scalar multiple of $\partial$.
To identify the scalar, use the canonical self-duality of an elliptic
curve to identify $H^1(E,\mathcal O_E)$ with the invariant tangent
space. Under this identification, Frobenius on $H^1(E,\mathcal O_E)$
acts by taking the $p$th iterate of an invariant derivation. Its
coefficient in the basis dual to $\omega$ is the definition of
$a_\omega$. This is the standard description in
\cite[\S2.0, pp.~97--98]{katz}, and proves
\eqref{eq:restricted-Hasse} with the stated normalization.
\end{proof}


\begin{thm}[Map-independent Hasse reduction]\label{thm:general-Hasse}
Let $R$ be a $p$-torsion-free $\mathbb Z_{(p)}$-algebra, let $E/R$ be
an elliptic scheme, and choose a frame $\omega$ of its invariant
differentials. Let $U\subset E$ be an affine open and $f\in\mathcal O(U)$
with $\partial f\in\mathcal O(U)^\times$, where $\partial$ is dual to
$\omega$. Then
\[
\mathcal S_{p-1}^\omega[f]
:=-\frac{(p-1)!}{B_{p-1}}S_{p-1}^\partial[f]
\in\mathcal O(U),
\]
and, on $\bar U$,
\begin{equation}\label{eq:general-Hasse-reduction}
\overline{\mathcal S_{p-1}^\omega[f]}=a_{\bar\omega}.
\end{equation}
Consequently the reduction of
$\mathcal S_{p-1}^\omega[f]\omega^{\otimes(p-1)}$
is the restriction of the pullback of $A$. On every geometric special
fiber meeting $\bar U$, it extends uniquely to the entire fiber as
that regular section. The extension is independent of $f$ and of its
integral lift, and vanishes identically precisely when the fiber is
supersingular.
\end{thm}

\begin{proof}
Smoothness of $E/R$ and $p$-torsion-freeness of $R$ imply
$p$-torsion-freeness of $\mathcal O(U)$. The invariant derivation
preserves this ring. Apply Proposition~\ref{prop:Bernoulli-reduction}
and then Lemma~\ref{lem:restricted-Hasse} to obtain
\[
\overline{\mathcal S_{p-1}^\omega[f]}
=\frac{\bar\partial^{\,p}\bar f}{\bar\partial\bar f}
=a_{\bar\omega}.
\]
The denominator is a unit by hypothesis. Multiplying by
$\bar\omega^{\otimes(p-1)}$ gives the asserted equality of sections.
A nonempty open subset of a geometric elliptic curve is dense, so this
equality determines its regular extension uniquely. The vanishing
criterion is the usual one for the Hasse invariant
\cite[\S2.0]{katz}.

For a unit $c$ pulled back from the base, replacing $\omega$ by
$c\omega$ replaces $\partial$ by $c^{-1}\partial$ and
$S_{p-1}^\partial[f]$ by $c^{-(p-1)}S_{p-1}^\partial[f]$.
The tensor therefore glues under changes of frame. Its reduction is
independent of the map and its lift because it is the pullback of $A$.
\end{proof}

\begin{thebibliography}{99}
\bibitem[ireland-rosen]{ireland-rosen} K.~Ireland and M.~Rosen, \emph{A Classical Introduction to Modern Number Theory}, second edition, Graduate Texts in Mathematics 84, Springer, New York, 1990.
\bibitem[katz]{katz} N.~M.~Katz, $p$-adic properties of modular schemes and modular forms, in \emph{Modular Functions of One Variable III}, Lecture Notes in Math. 350, Springer, Berlin, 1973, 69--190.
\end{thebibliography}
\end{document}
